% GENERATED FILE. Edit canonical lessons, then run npm run generate:books.
% canonical-source-hash: fnv1a64:c330c90fde8fb5d6
% canonical-lessons: FR-C375-lycee, FR-C375-lyceen, FR-C375-frequenter, FR-C375-bulletin, FR-C375-proviseur

\chapter{Au lycée}
\label{ch:fr-au-lycee}
\hlchaptermodality{375}

\begin{chapteropening}
\textbf{By the end of this chapter:} \emph{I can talk about secondary school: the lycée, its pupils, attending it, the school report and the head.}

Complete the last lesson of chapter 375: I can talk about secondary school: the lycée, its pupils, attending it, the school report and the head.
\end{chapteropening}

\section[le lycée]{le lycée — the secondary school (ages fifteen to eighteen)}

\label{lesson:FR-C375-lycee}



\begin{quote}
Before the new one: say the French for the resolution, then the French for to aim.
\end{quote}

\subsection*{You'll want to know: le lycée}
\textbf{le lycée} — \textquotedblleft{}the secondary school (ages fifteen to eighteen)\textquotedblright{}.

\begin{quote}
Ma sœur entre au lycée en septembre.
\end{quote}

\begin{quote}
\emph{My sister starts secondary school in September.}
\end{quote}

Belgium: \emph{l'athénée} or \emph{l'école secondaire}; Québec: \emph{l'école secondaire}, then \emph{le cégep}; Switzerland: \emph{le gymnase}.

\begin{cousinweb}
\emph{lycée} comes from Latin \emph{lycēum} (EtymDB), from Greek \emph{Lúkeion}, the name of a gymnasium where a famous philosopher taught.
\end{cousinweb}

\begin{grammarlens}[title={What you've built}]
One more word for telling, explaining and describing at B1.
\end{grammarlens}

\subsection*{Your turn}
\begin{itemize}
\raggedright
  \item \emph{Say it:} \emph{le lycée}
  \item \emph{Say it:} \emph{le lycée}, once more
  \item \emph{Say it:} a sentence of your own with \emph{le lycée}
  \item \emph{Recall:} say the French for the resolution, then the French for to aim, then say \emph{le lycée} again
\end{itemize}

\begin{tcolorbox}[breakable,colback=teal!4,colframe=teal!35!black,title={Before you move on}]
What does le lycée mean? (\textquotedblleft{}The secondary school (ages fifteen to eighteen)\textquotedblright{}.) Say it once more.
\end{tcolorbox}
\section[le lycéen]{le lycéen — the secondary-school pupil}

\label{lesson:FR-C375-lyceen}



\begin{quote}
Before the new one: say the French for to aim, then the French for the secondary school.
\end{quote}

\subsection*{You'll want to know: le lycéen}
\textbf{le lycéen} — \textquotedblleft{}the secondary-school pupil\textquotedblright{}.

\begin{quote}
Les lycéens ont des cours jusqu'à six heures.
\end{quote}

\begin{quote}
\emph{The pupils have lessons until six o'clock.}
\end{quote}

Feminine \emph{lycéenne}.

\begin{cousinweb}
\emph{lycéen} is built from \emph{lycée} and \emph{-en}. \emph{lycée} comes from Latin \emph{lycēum}, from Greek \emph{Lúkeion} (EtymDB).
\end{cousinweb}

\begin{grammarlens}[title={What you've built}]
One more word for telling, explaining and describing at B1.
\end{grammarlens}

\subsection*{Your turn}
\begin{itemize}
\raggedright
  \item \emph{Say it:} \emph{le lycéen}
  \item \emph{Say it:} \emph{le lycéen}, once more
  \item \emph{Say it:} a sentence of your own with \emph{le lycéen}
  \item \emph{Recall:} say the French for to aim, then the French for the secondary school, then say \emph{le lycéen} again
\end{itemize}

\begin{tcolorbox}[breakable,colback=teal!4,colframe=teal!35!black,title={Before you move on}]
What does le lycéen mean? (\textquotedblleft{}The secondary-school pupil\textquotedblright{}.) Say it once more.
\end{tcolorbox}
\section[fréquenter]{fréquenter — to attend, to go regularly to; to spend time with}

\label{lesson:FR-C375-frequenter}



\begin{quote}
Before the new one: say the French for the secondary school, then the French for the secondary-school pupil.
\end{quote}

\subsection*{You'll want to know: fréquenter}
\textbf{fréquenter} — \textquotedblleft{}to attend, to go regularly to; to spend time with\textquotedblright{}.

\begin{quote}
Il fréquente un lycée près de chez lui.
\end{quote}

\begin{quote}
\emph{He attends a secondary school near his home.}
\end{quote}

\emph{Fréquenter quelqu'un} — to see a lot of someone.

\begin{cousinweb}
\emph{fréquenter} is a learned borrowing from Latin \emph{frequentāre} \textquotedblleft{}to visit often, to crowd\textquotedblright{} (Lewis \& Short), from \emph{frequēns} \textquotedblleft{}crowded, frequent\textquotedblright{}.
\end{cousinweb}

\begin{grammarlens}[title={What you've built}]
One more word for telling, explaining and describing at B1.
\end{grammarlens}

\subsection*{Your turn}
\begin{itemize}
\raggedright
  \item \emph{Say it:} \emph{fréquenter}
  \item \emph{Say it:} \emph{fréquenter}, once more
  \item \emph{Say it:} a sentence of your own with \emph{fréquenter}
  \item \emph{Recall:} say the French for the secondary school, then the French for the secondary-school pupil, then say \emph{fréquenter} again
\end{itemize}

\begin{tcolorbox}[breakable,colback=teal!4,colframe=teal!35!black,title={Before you move on}]
What does fréquenter mean? (\textquotedblleft{}To attend, to go regularly to; to spend time with\textquotedblright{}.) Say it once more.
\end{tcolorbox}
\section[le bulletin]{le bulletin — the school report; the bulletin}

\label{lesson:FR-C375-bulletin}



\begin{quote}
Before the new one: say the French for the secondary-school pupil, then the French for to attend.
\end{quote}

\subsection*{You'll want to know: le bulletin}
\textbf{le bulletin} — \textquotedblleft{}the school report; the bulletin\textquotedblright{}.

\begin{quote}
Cette année, mon bulletin est meilleur.
\end{quote}

\begin{quote}
\emph{This year, my school report is better.}
\end{quote}

\emph{Un bulletin de vote} — a ballot paper.

\begin{cousinweb}
\emph{bulletin} was borrowed from Italian \emph{bullettino} \textquotedblleft{}a small note, a ticket\textquotedblright{} (EtymDB), built on Latin \emph{bulla} \textquotedblleft{}a bubble, a seal\textquotedblright{} (Lewis \& Short).
\end{cousinweb}

\begin{grammarlens}[title={What you've built}]
One more word for telling, explaining and describing at B1.
\end{grammarlens}

\subsection*{Your turn}
\begin{itemize}
\raggedright
  \item \emph{Say it:} \emph{le bulletin}
  \item \emph{Say it:} \emph{le bulletin}, once more
  \item \emph{Say it:} a sentence of your own with \emph{le bulletin}
  \item \emph{Recall:} say the French for the secondary-school pupil, then the French for to attend, then say \emph{le bulletin} again
\end{itemize}

\begin{tcolorbox}[breakable,colback=teal!4,colframe=teal!35!black,title={Before you move on}]
What does le bulletin mean? (\textquotedblleft{}The school report; the bulletin\textquotedblright{}.) Say it once more.
\end{tcolorbox}
\section[le proviseur]{le proviseur — the head teacher (of a lycée)}

\label{lesson:FR-C375-proviseur}



\begin{quote}
Before the new one: say the French for to attend, then the French for the school report.
\end{quote}

\subsection*{You'll want to know: le proviseur}
\textbf{le proviseur} — \textquotedblleft{}the head teacher (of a lycée)\textquotedblright{}.

\begin{quote}
Le proviseur a parlé aux parents hier soir.
\end{quote}

\begin{quote}
\emph{The head teacher spoke to the parents last night.}
\end{quote}

In a \emph{collège}, the head is \emph{le principal}. Belgium: \emph{le préfet}; Québec: \emph{le directeur}.

\begin{cousinweb}
\emph{proviseur} is a learned borrowing from Latin \emph{prōvīsor} \textquotedblleft{}one who foresees, a provider\textquotedblright{} (EtymDB; Lewis \& Short), from \emph{prōvidēre} \textquotedblleft{}to foresee, to provide for\textquotedblright{}.
\end{cousinweb}

\begin{grammarlens}[title={What you've built}]
One more word for telling, explaining and describing at B1.
\end{grammarlens}

\subsection*{Your turn}
\begin{itemize}
\raggedright
  \item \emph{Say it:} \emph{le proviseur}
  \item \emph{Say it:} \emph{le proviseur}, once more
  \item \emph{Say it:} a sentence of your own with \emph{le proviseur}
  \item \emph{Recall:} say the French for to attend, then the French for the school report, then say \emph{le proviseur} again
\end{itemize}

\begin{tcolorbox}[breakable,colback=teal!4,colframe=teal!35!black,title={Before you move on}]
What does le proviseur mean? (\textquotedblleft{}The head teacher (of a lycée)\textquotedblright{}.) Say it once more.
\end{tcolorbox}
